% Independent all-weight proof; baseline v32 is unchanged.
\section{任意の重みで増大する一般点の非被約厚み}
\label{sec:all-weights-concurrent-v33}

この節では整数 $m\geq3$ を固定し、
$X_m=\mathbb P_{\mathrm{stack}}(1,1,1,m)$ 上の大域的 Poisson テンソルの
アフィンスキームを $\mathscr P_m$ と書く。係数体は標数零とする。
$S_d=k[x,y,z]_d$、$N_m=\dim S_m=\binom{m+2}{2}$ と置く。

方程式と記号は命題~\ref{prop:general-equations-v33}のものを用いる。
最初に、$m$ に依存しない基底切片を決定する。
% Complete base slice before imposing the mixed F equations.
\subsection{共点三直線における基底の非被約な形式切片}
\label{subsec:concurrent-base-v32}

$\mathscr Z=\{(c,D)\in S_3\oplus\mathfrak{sl}_3:D(c)=0\}$ と置く。
この基底スキーム自身に既に非被約構造が現れる。
以下の補題は混合項 $F$ の方程式をまだ課さない、基底だけの全形式的な記述である。

\begin{lemma}[共点三直線の基底切片]
\label{lem:concurrent-base-v32}
\[
 c_0=xy(x+y),\qquad D_0=(Ax+By)\partial_z
\]
とし、
\[
 U=B-2A,\qquad V=A-2B,\qquad W=A+B
\]
が全て非零とする。例えば $(A,B)=(2,5)$ はこの条件を満たす。
この点における $\mathscr Z$ の完備局所環は
\[
 \frac{\mathbb C[[u_1,\ldots,u_8,a,b]]}{(ab,b^2)}
\]
に同型である。八つの滑らかな変数のうち六つは $c_0$ の
$\mathrm{SL}_3$ 軌道方向、残る二つは $A,B$ の変動である。
被約化は滑らかな九次元であり、零でない平方零イデアル $(b)$ は
共点三直線の部分 $a=0$ に支持される。
\end{lemma}

\begin{proof}
$c_0$ の $\mathfrak{sl}_3$ による接軌道は六次元である。
その像は二変数三次式の四次元空間と
$z\partial_xc_0,z\partial_yc_0$ の二方向からなる。
従って、四次元の形式横断切片を
\[
 c=xy(x+y)+a xyz+b xz^2+\gamma yz^2+d z^3
\]
と取れる。六つの独立な軌道方向を形式的な群座標として選べば、
軌道方向とこの切片から $S_3$ への写像の接写像は同型であり、形式逆写像定理が適用できる。
$D$ も同じ $\mathrm{SL}_3$ 座標変換で移すので、$D(c)=0$ の問題は
この六次元の滑らかな因子と切片上の問題に分かれる。

切片上で $D_{zx}=A,D_{zy}=B$ と書く。
$D(c)=0$ の四つの二変数三次単項式と $x^2z,y^2z$ の係数から、
残る六つの係数は一意に
\[
\begin{aligned}
 D_{yx}=D_{xy}&=0,\\
 D_{xx}&=\frac{a(B-2A)}3,&
 D_{yy}&=\frac{a(A-2B)}3,\\
 D_{xz}&=-2B\gamma,&D_{yz}&=-2Ab,\\
 D_{zz}&=-D_{xx}-D_{yy}
\end{aligned}
\]
と決まる。代入後の残る四式は、零でない定数倍を除いて
\[
\begin{aligned}
 E_1&=Ub+V\gamma,\\
 E_2&=aUb+3Ad,\\
 E_3&=aV\gamma+3Bd,\\
 E_4&=W(ad-2b\gamma)
\end{aligned}
\]
である。

この段階で $U,V,W$ は形式局所環の単元である。まず
\[
 E_2+E_3-aE_1=3Wd
\]
から $d$ を消去でき、次に $E_2-3Ad=Uab$ から $ab=0$ を得る。
また $E_1=0$ により $\gamma=-Ub/V$ である。
最後に $E_4=0$ から $b\gamma=0$、従って $b^2=0$ が従う。
逆に $d=0,\gamma=-Ub/V,ab=b^2=0$ を満たせば四つの $E_i$ は全て零になる。
従ってこれは零点集合だけの比較でなく、局所化された関係イデアルの等式である。

残った自由変数は六つの軌道座標と $A,B,a,b$ であり、表示した完備局所環を得る。
$\mathbb C[[a,b]]/(ab,b^2)$ において $b\ne0$、$b^2=ab=0$ だから、
非被約性とその支持も明らかである。
なお、消去には $A$ や $B$ 自身の可逆性は必要なく、上記三つの単元条件で十分である。
\end{proof}

被約化の切片では $c=xy(x+y+az)$ となる。
従って $a\ne0$ は非共点の三直線、$a=0$ は共点三直線に対応する。
この補題だけから、混合項を含む全 Poisson スキームの既約成分性や
一般点の非被約厚みの次数を結論してはならない。
その判定には、この環上での混合写像をさらに解析する必要がある。


\subsection{全重みに共通する局所環}

\begin{theorem}[一般点の厚み $m+3$]
\label{thm:allm-concurrent-v33}
各整数 $m\geq3$ に対し、$(A,B)=(2,5)$ を含む非空 Zariski 開集合が存在し、
\[
 c_0=xy(x+y),\qquad D_0=(Ax+By)\partial_z,\qquad [F]=0
\]
における $\mathscr P_m$ の完備局所環は
\[
 \boxed{\quad
 \frac{k[[s_1,\ldots,s_{N_m+8},a,b,u]]}
 {(ab,b^2,a^{m+2}u)}
 \quad}
\]
と同型である。開集合は $m$ に依存してよい。
被約化の二つの枝の次元はともに $N_m+9$、その被約交叉は滑らかな
$N_m+8$ 次元である。$u\ne0$ の枝の一般点では横断局所環が
\[
 \frac{K[[a,b]]}{(a^{m+2},ab,b^2)}
\]
となり、その長さは $m+3$ である。
ここで $K$ は当該被約成分の関数体であり、係数体の選択は非標準的である。
\end{theorem}

以下の証明は任意の $m$ を一度に扱う。
有限個の重みにおける行列計算から帰納的な一般化を推測するものではない。

\subsubsection*{基底切片と二つの次数上昇}

$D(c)=0$ は $m$ に依存しない。
共点三直線の基底切片の消去により、六つの軌道座標を別にして
\begin{align*}
 c&=xy(x+y)+a xyz+b(xz^2+\gamma yz^2),\\
 D_x&=\frac{a(B-2A)}3x-2B\gamma bz,\\
 D_y&=\frac{a(A-2B)}3y-2Abz,\\
 D_z&=Ax+By+\frac{a(A+B)}3z,\qquad
 \gamma=-\frac{B-2A}{A-2B},
\end{align*}
かつ関係は正確に $ab=b^2=0$ となる。
この消去には $(B-2A)(A-2B)(A+B)\ne0$ だけを要し、$(2,5)$ で成立する。
$A,B$ の二変数と合わせて八つの滑らかな基底変数を得る。

混合写像を $M=M_0+aM_1+bM_2$ と書く。
$F_i=x^{n_1}y^{n_2}z^{n_3}$ の重量を
$n_3-\mathbf1_{i=3}$ とし、標的には $z$ 次数を与える。
このとき $M_0,M_1,M_2$ はそれぞれ重量を $0,1,2$ 上げる。
この主張は各項の次数から直接従い、係数 $A,B$ や $m$ には依存しない。

\subsubsection*{共点側の核を全重みで計算する}

$L=Ax+By$ と置き、二変数の二次ベクトル場
\[
 V=\frac13\bigl((c_{0,y}+xL)\partial_x
              -(c_{0,x}-yL)\partial_y\bigr)
\]
を取る。$\Pi_0=\partial_z\wedge V$ であり、
$V(c_0)=Lc_0$、$\operatorname{div}V=L$ である。
$F=Q+H\partial_z$ と書く。
混合式を三直線へ制限すると、次の三つの係数が非零である限り $Q$ は各直線に接する。
\[
 \kappa_1=1+\frac m3(1-B),\qquad
 \kappa_2=1+\frac m3(1+A),\qquad
 \kappa_3=1+\frac m3(1-A+B).
\]
実際、$x=0$ では $\kappa_1 y^2F_x=0$、$y=0$ では
$\kappa_2 x^2F_y=0$、$x+y=0$ では
$-\kappa_3x^2(F_x+F_y)=0$ となる。
この条件は $(2,5)$ において
\[
 (\kappa_1,\kappa_2,\kappa_3)
 =\left(1-\frac{4m}3,1+m,1+\frac{4m}3\right)
\]
なので、全ての $m\geq3$ で成立する。

三直線に接する二変数導分の自由基底は $E_2,V$ である。
ここで $\det(E_2,V)=-c_0$ であるから、これは係数 $A,B$ によらない。
Euler 同値で $E_2$ 成分を消すと、重量 $j\geq0$ の核は
\[
 Q=z^jgV,\qquad H=-\frac{z^{j+1}}{j+1}V(g),
 \qquad g\in k[x,y]_{m-1-j},\quad 0\leq j\leq m-1.
\]
なぜなら混合式の左辺は、この代表に対して
$-c_0(V(g)z^j+\partial_zH)$ となるからである。
これらは全て $-z^{j+1}g/(j+1)$ による次数 $m$ の $w$ 移動である。
重量 $-1$ では $H=f(x,y)$、$\deg f=m+1$ の $m+2$ 方向が残る。
従って
\[
 \dim\ker M_0=\frac{m(m+1)}2+(m+2)=N_m+1.
\]

二変数の次数 $m$ の $w$ 移動を明示すると、
$g_k=x^{m-k}y^k$、$e_k=x^{m+1-k}y^k$ に対し
\[
 V(g_k)=\left(\frac{m(A+1)}3-k\right)e_k
       +\left(\frac{m(B+2)}3-k\right)e_{k+1},
 \quad 0\leq k\leq m.
\]
$(2,5)$ では後者の係数が $7m/3-k>0$ である。
従って二変数移動の像は $m+1$ 次元、その余核は $e_0=x^{m+1}$ で生成される。
非二変数の移動の単射性は上の $Q=z^jgV$ の表示から分かる。
結局、全ての次数 $m$ の移動の像は $N_m$ 次元で、混合核との商は
重量 $-1$ の一点
\[
 F_{\mathrm{ex}}=(0,0,x^{m+1})
\]
で生成される。

以上から重量ブロックの大きさと階数は正確に次のようになる。
\[
\begin{array}{c|c|c|c}
 \text{重量 }r&\text{源}&\text{標的}&\operatorname{rank}M_0\\\hline
 -1&m+2&0&0\\
 0\leq r\leq m-1&2(m+2-r)&m+4-r&m+4-r\\
 m&4&4&4\\
 m+1&2&3&2\\
 m+2&0&2&0\\
 m+3&0&1&0
\end{array}
\]
源の次元は、例えば $F_x$ の $x$ で割り切れる単項式を Euler 同値で
消すことにより数えられる。余核は四次元で、重量は
$m+1,m+2,m+2,m+3$ である。

\subsubsection*{非共点側の核と非共鳴条件}

$b=0$ とし、$a$ を反転させる。
$X=x,Y=y,Z=x+y+az$ の線形座標変換を行うと、
係数ベクトルの体積変換を含めて
\[
 c'=aXYZ,\qquad
 D'=\frac a3\bigl((B-2A)X\partial_X
 +(A-2B)Y\partial_Y+(A+B)Z\partial_Z\bigr).
\]
標的の単項式 $X^{\nu_1}Y^{\nu_2}Z^{\nu_3}$、
$\sum\nu_i=m+3$ を固定する。
全ての $\nu_i$ が正なら、対応する源は Euler 同値の前で三次元、
同値の後で二次元であり、$M/a$ の行は
\[
 (\kappa_1-\nu_1,\kappa_2-\nu_2,\kappa_3-\nu_3).
\]
$\sum\kappa_i=m+3$ なので、この行は Euler 同値を殺す。
その階数が零となるのは $(\kappa_1,\kappa_2,\kappa_3)=\nu$ の場合だけである。
一つの $\nu_i$ が零で、他の二つが正なら、源は一次元で係数は $\kappa_i$ である。
二つ以上が零なら源は存在しない。

従って、$\kappa_i\ne0$ かつ $\kappa$ が正整数三つからなる上記有限集合を
避けるとき、核の次元は内部単項式の数 $N_m$ に等しい。
$(2,5)$ では $\kappa_1<0$ なので、この条件も全ての $m\geq3$ で成立する。
次数 $m$ の $w$ 移動は常に実際の核に入るから、ここではその像が核の全体である。
特に $a$ を反転すると、混合写像の階数は共点側から正確に一つ増える。

\subsubsection*{Schur 補完と三交点評価}

各重量で $M_0$ の可逆小行列を選び、対応する源の変数を消去する。
$r\leq m$ の標的行は全て消去し、$r=m+1$ では
$xy z^{m+1},y^2z^{m+1}$ の二行を消去する。
後者の二次小行列式は $(2,5)$ において
\[
 -\frac{(3-4m)(6+7m)}9\ne0
\]
である。従って、消去後に残る標的の四行は
\[
 x^2z^{m+1},\quad xz^{m+2},\quad yz^{m+2},\quad z^{m+3}
\]
と取れる。必要な有限個の小行列式を反転させた非空開集合上で、この選択は保たれる。

次数 $m$ の移動の写像
\[
 T:S_m\longrightarrow S_{m+1}^3/S_m\mathbf x,
 \qquad g\longmapsto-\,p\times\nabla g
\]
は $D(c)=0$ 上で常に混合方程式を満たす。
特殊点における階数が $N_m$ なので、$N_m$ 次小行列式が局所環の単元であり、
$T(S_m)$ は非被約基底 $k[[a,b]]/(ab,b^2)$ 上でも自由な直和因子である。
Schur 消去後にも、その像は残った源の自由加群の直和因子となる。
実際、Schur 消去で消した座標は残りの座標から一意に復元されるから、
特殊点における像の階数は変わらない。
従って $N_m$ 個の核座標と、残余の一座標 $u$ を選べる。

$b=0$ とした、まだ核座標を取り直していない Schur 写像を $R(a)$ とする。
可逆小行列の定数部は重量を保ち、その $a$ 部は重量を一つ上げる。
従って逆行列の幾何級数は重量の有限性により有限和である。
重量 $w$ の自由列から重量 $r$ の自由行への成分は、零でなければ
正確に $a^{r-w}$ の定数倍である。
特に $r=m+1$ の行では、全ての成分の $a$ 次数は $m+2$ 以下である。

$R(a)$ の最初の非零 $a$ 係数は
$\ker M_0/T(S_m)$ を通る。
これは $R(a)T_{\mathrm{free}}(a)=0$ の最初の非零係数を取り出せば分かる。
この商の重量は $-1$、余核の最小重量は $m+1$ である。
従って残余写像は $a^{m+2}$ で割り切れる。
同じ議論で $b$ の一次係数は重量を二つしか上げず、余核に届かないため零である。
$ab=b^2=0$ なので、$b$ を含む他の係数もない。

ここで先頭 $a^{m+2}$ 係数の非零性を、全ての $m$ に対して証明する。
Schur 消去の直後、残余出力は上記四行だけに支持される。
$b=0$ では $c=xy(x+y+az)$ であり、三つの交点
\[
 [0,0,1],\qquad[-a,0,1],\qquad[0,-a,1]
\]
において混合写像の像は消える。
$a\ne0$ では三角形座標の単項式表示から明らかであり、
多項式恒等式として $a=0$ にも及ぶ。
出力を
\[
 r_0x^2z^{m+1}+r_1xz^{m+2}+r_2yz^{m+2}+r_3z^{m+3}
\]
と書けば、三交点での評価から
\[
 r_3=0,\qquad a^2r_0-ar_1=0,\qquad-ar_2=0.
\]
$b=0$ の係数環は整域なので、$r_1=ar_0,r_2=r_3=0$ である。
次数上昇と残余商の重量から、$r_0$ の最初に可能な項は
$a^{m+2}$ である。一方、まだ核座標を取り直していない $R(a)$ の
$r_0$ 行は、前段により $a$ の次数 $m+2$ 以下の多項式である。
従ってこの係数が零なら $r_0$ 行は恒等的に零となる。
従って先頭係数が零なら残余写像全体が零となるが、これは $a$ を反転したとき
階数が一つ増えることに矛盾する。

従って $(2,5)$ で先頭係数は非零であり、非空開集合を縮めれば単元である。
残余の関係イデアルは正確に $(a^{m+2}u)$ となる。
八つの滑らかな基底座標と $N_m$ 個の移動座標を合わせて、定理の局所環を得る。

\subsubsection*{一般点の長さ}

被約化の二つの素イデアルは $(a,b)$ と $(b,u)$ である。
前者の枝で $u$ を反転すると、横断環は
\[
 K[a,b]/(a^{m+2},ab,b^2)
\]
となり、その $K$ 基底は $1,a,\ldots,a^{m+1},b$ である。
係数体を幾何学的な拡大へ取る必要はない。
完備局所環の係数体を選び、長さ $m+2$ の主鎖を与える $a$ と、
その最終項から独立な socle 元 $b$ を選べば、表示した $K$ 形式を得る。
これで定理が証明された。

\begin{remark}[計算による確認の役割]
付属の有理数検証コードは複数の重みで重量ブロック、Schur 消去、三交点評価、
$b$ 項の消滅を独立に確認する。
全重みの結論は、ここで証明した核の分解、三角形側の非共鳴条件、
重量と三交点評価から従うのであり、有限標本に基づく推定ではない。
\end{remark}
